\section{Тема 12. Кривые второго порядка на плоскости}

\subsection{Постановка задачи (вариант~4)}

Из круга радиуса $R$ вырезан прямоугольник со сторонами $w$ и $h$,
расположенный произвольно: с центром $(c_x,c_y)$ и углом
поворота~$\theta$ относительно координатных осей.
Найти координаты его четырёх вершин.

\subsection{Математическое обоснование}

В локальной (неповёрнутой) системе координат вершины прямоугольника
с центром в начале координат:
\[
  (\pm w/2,\,\pm h/2).
\]
Поворот точки $(dx,dy)$ на угол~$\theta$ выполняется матрицей
поворота:
\[
  \begin{pmatrix}x'\\y'\end{pmatrix} =
  \begin{pmatrix}\cos\theta & -\sin\theta\\ \sin\theta & \cos\theta\end{pmatrix}
  \begin{pmatrix}dx\\dy\end{pmatrix}.
\]
Итоговые координаты вершины:
\[
  (c_x + x',\; c_y + y').
\]

\subsection{Блок-схема алгоритма}

\begin{center}
\begin{tikzpicture}[
  node distance=10mm, start chain=going below,
  nodes={draw=black, top color=white, bottom color=lightgray!50,
         minimum width=5.6cm, align=center, font=\small},
  arrow/.style={->, >=Stealth, thick},
]
  \node[draw, rounded corners, on chain] {Start};
  \node[shape=trapezium, trapezium left angle=70,
        trapezium right angle=110, on chain, draw, join=by arrow]
       {$c_x,c_y,\ w,h,\ \theta$};
  \node[shape=rectangle, on chain, draw, join=by arrow]
       {$dx[4],\ dy[4]$ -- локальные вершины};
  \node[shape=chamfered rectangle, chamfered rectangle xsep=10pt,
        on chain, draw, join=by arrow]
       (loop) {$i := 0,\ 3$};
  \begin{scope}[local bounding box=scl]
    \node[shape=rectangle, on chain, draw, join=by arrow]
         {$(x',y') :=$ поворот $(dx_i,dy_i)$ на $\theta$};
    \node[shape=trapezium, trapezium left angle=70,
          trapezium right angle=110, on chain, draw, join=by arrow]
         {$(c_x{+}x',\ c_y{+}y')$};
    \coordinate[on chain, join=by arrow] (le);
  \end{scope}
  \draw[arrow] (scl.south) -| ($(scl.west)-(0.9,0)$) |- (loop);
  \coordinate[on chain] (sp);
  \draw[arrow] (loop) -| ($(scl.east)+(0.9,0)$) |- (sp);
  \node[draw, rounded corners, on chain, join=by arrow] {End};
\end{tikzpicture}
\end{center}

\subsection{Текст программы}

\lstinputlisting[style=Cstyle]{../src/t12.c}

\subsection{Тестовые данные}

\begin{center}
\begin{tabular}{l|l}
\toprule
Параметры & Вершины \\
\midrule
$c=(0,0),\,w=2,h=2,\theta=0°$ & $(1,1),(-1,1),(-1,-1),(1,-1)$ \\
\bottomrule
\end{tabular}
\end{center}

\textbf{Результат выполнения программы} (центр $(5,5)$, $2{\times}3$, поворот $45°$):
\begin{lstlisting}[language={},basicstyle=\ttfamily\small,frame=single]
Центр прямоугольника (cx cy): 5 5
Ширина и высота (w h): 2 3
Угол поворота (градусы): 45
Вершины прямоугольника:
  (5.3536, 3.2322)
  (6.7678, 4.6464)
  (4.6464, 6.7678)
  (3.2322, 5.3536)
\end{lstlisting}
